\textbf{Codex's integration has advanced: PR \#16 is merged, and the combined frontier/Codex run has
now passed.} I checked the current repository state and the completed job: it rebuilt the inherited
sources, checked the frontier files, imported both developments into one Lean environment, and
completed their selected axiom audits. The ``combined run pending'' sentence in PR \#18's
description is therefore out of date.

I continued with the new sum-connection note rather than duplicating that work. \textbf{The
additional result is a complete calculation of what its conormal derivative detects, what it misses,
and how that calculation restricts to the symmetric tensor summand.} It also packages the quotient
and derivative together as a genuine algebra homomorphism.

The new results below have complete mathematical proofs and exact finite algebra checks. The Lean
continuation at the end is an \textbf{uncompiled draft}: the local execution service is timing out,
and the currently exposed GitHub actions provide no write operation. No remote files were changed in
this turn.

\subsubsection{1. The derivative now has a complete quotient
construction}\label{conormal_tower_pasted:the-derivative-now-has-a-complete-quotient-construction}

Retain the source's notation, with tensor degree \(k\ge1\):

\[
P=\mathbb C[s_1,\ldots,s_k],\qquad
h_i=h(s_i),\qquad
I=(h_1,\ldots,h_k),
\]

\[
E=P/I,\qquad E^{[2]}=P/I^2,
\qquad
\partial_S=\frac1k\sum_{i=1}^k\partial_{s_i}.
\]

Here \(h\) is the original monic packet polynomial, including all multiplicities. These are
precisely the relation ideal and derivative in the new attachment.

For any derivation \(D\) of a commutative algebra and any ideal \(I\),

\[
\boxed{D(I^{r+1})\subseteq I^r\qquad(r\ge0).}
\]

\textbf{Proof.} An element of \(I^{r+1}\) is a sum of multiples of products of \(r+1\) elements of
\(I\). The product rule differentiates one factor at a time, leaving at least \(r\) factors in
\(I\). Differentiating an exterior coefficient leaves all \(r+1\) factors. Additivity completes the
proof.

Consequently the source's maps exist at every order:

\[
\delta_r:P/I^{r+1}\longrightarrow P/I^r,
\qquad
[F]\longmapsto[\partial_SF].
\]

They commute with the quotient transitions. In particular,

\[
\delta_1:E^{[2]}\longrightarrow E
\]

is a derivation \textbf{along} the quotient \(\pi:E^{[2]}\to E\):

\[
\boxed{
\delta_1(xy)=\pi(x)\delta_1(y)+\pi(y)\delta_1(x).
}
\]

This is the appropriate product rule; \(\delta_1\) is not an endomorphism of the arithmetic algebra
\(E\).

\paragraph{Monic division constructs all the retained
layers}\label{conormal_tower_pasted:monic-division-constructs-all-the-retained-layers}

There is a unique finite expansion

\[
F=\sum_{\alpha\in\mathbb N^k}
h_1^{\alpha_1}\cdots h_k^{\alpha_k}r_\alpha,
\qquad
\deg_{s_i}r_\alpha<d:=\deg h.
\]

Existence follows by repeated monic division in the separate variables. For uniqueness, the leading
monomials of

\[
h^\alpha s^\beta,\qquad 0\le\beta_i<d,
\]

have exponent vectors \(d\alpha+\beta\), which are distinct. A finite dependence therefore has a
unique largest leading monomial and cannot cancel.

This gives explicit coordinates---not just an existence statement---for

\[
I^r/I^{r+1}\cong E[\eta_1,\ldots,\eta_k]_r,
\qquad
\eta_i\longleftrightarrow[h_i].
\]

In particular,

\[
I/I^2\cong E^k,
\qquad
(a_i)\longmapsto\left[\sum_i h_i\widetilde a_i\right],
\]

and

\[
\dim_{\mathbb C}P/I^{r+1}
=d^k\binom{k+r}{k}.
\]

On the associated graded, the derivative is exactly

\[
\boxed{
\operatorname{gr}(\partial_S)
=
\frac1k\sum_i[h_i']\,\partial_{\eta_i}.
}
\]

Indeed, differentiating a coefficient of \(h^\alpha r_\alpha\) stays in \(I^{|\alpha|}\), while
differentiating one \(h_i\)-factor lowers the ideal order and contributes the literal exponent
\(\alpha_i\). This proves the higher-layer formula in the attachment without deleting those exponent
factors.

The standard conormal construction supplies the map

\[
I/I^2\longrightarrow\Omega_{P/\mathbb C}\otimes_PE.
\]

The calculation here identifies its coordinates for this particular ideal and then contracts it by
the specified sum direction; those are distinct operations (The Stacks Project Authors, n.d., Tag
00S0). (\href{https://stacks.math.columbia.edu/tag/00S0}{The Stacks Project})

\subsubsection{2. The exact information detected by the sum
derivative}\label{conormal_tower_pasted:the-exact-information-detected-by-the-sum-derivative}

Let

\[
A_h=\mathbb C[s]/(h),\qquad E=A_h^{\otimes k},
\]

and write

\[
c_i=[h'(s_i)]\in E.
\]

Under the preceding conormal coordinates, the restriction of \(\delta_1\) is the \(E\)-linear map

\[
\boxed{
\partial_N:E^k\longrightarrow E,
\qquad
(a_1,\ldots,a_k)\longmapsto\frac1k\sum_i c_i a_i.
}
\]

The attachment gives this row. The following kernel, cokernel, and rank calculations are additional
deductions.

Let

\[
b=\deg\gcd(h,h').
\]

Over \(\mathbb C\), if \(h\) has \(q\) distinct roots, then

\[
b=d-q=\sum_\rho(m_\rho-1).
\]

The image and cokernel are

\[
\boxed{
\operatorname{im}\partial_N=(c_1,\ldots,c_k)\subseteq E,
}
\]

\[
\boxed{
\operatorname{coker}\partial_N
\cong
\left(\mathbb C[s]/(h,h')\right)^{\otimes k}.
}
\]

\textbf{Proof.} The image is the ideal generated by the row entries because \(1/k\) is invertible.
Quotienting \(E\) by those entries imposes \(h(s_i)=h'(s_i)=0\) in each separate variable, giving
the displayed tensor product.

Therefore

\[
\boxed{
\operatorname{rank}_{\mathbb C}\partial_N=d^k-b^k,
\qquad
\dim_{\mathbb C}\ker\partial_N=(k-1)d^k+b^k.
}
\]

These formulas separate two sources of information lost by this particular observation.

When \(h\) is squarefree, \(b=0\), so the contracted derivative is onto. Nevertheless, for \(k>1\)
its kernel still has dimension \((k-1)d^k\): \textbf{contracting to the sum direction does not
observe every relative conormal direction, even without repeated roots.}

When roots are repeated, the additional term \(b^k\) records the failure of the Jacobian entries to
generate the unit ideal.

For example, take the declared algebraic calibration

\[
h(s)=s^2,\qquad k=2,\qquad
E=\mathbb C[x,y]/(x^2,y^2).
\]

Then

\[
\partial_N(a,b)=xa+yb.
\]

Its image is

\[
\operatorname{span}\{x,y,xy\},
\]

so its rank is three and its kernel has dimension five. The conormal vector \((x,0)\) is nonzero but
lies in that kernel. Its representative is \(x^3\) modulo

\[
I^2=(x^4,x^2y^2,y^4),
\]

so the retained relation is genuinely nonzero before observation.

These are ranks of the \textbf{conormal restriction}, not of the whole map
\(\delta_1:E^{[2]}\to E\). The latter is onto: in coordinates

\[
S=\sum_i s_i,\qquad z_i=s_i-S/k,
\]

one has \(\partial_S=\partial/\partial S\), and a polynomial primitive is obtained by integrating
its powers of \(S\). The same argument makes every \(\delta_r\) onto.

\paragraph{Package the quotient and derivative as one algebra
map}\label{conormal_tower_pasted:package-the-quotient-and-derivative-as-one-algebra-map}

Introduce an ordinary formal square-zero parameter \(\epsilon\), distinct from the supported scalar
\(e\). Define

\[
\boxed{
\mathcal J_S^1:E^{[2]}
\longrightarrow E[\epsilon]/(\epsilon^2),
\qquad
x\longmapsto\pi(x)+\epsilon\,\delta_1(x).
}
\]

This is a unital \(\mathbb C\)-algebra homomorphism. Multiplication follows from the relative
product rule:

\[
(\pi x+\epsilon\delta x)(\pi y+\epsilon\delta y)
=
\pi(xy)+\epsilon\delta(xy).
\]

Thus the source's two differently typed maps can be retained together in a single multiplicative
observation. Its kernel is exactly

\[
\ker\mathcal J_S^1
=
\{n\in I/I^2:\partial_Nn=0\},
\]

and consequently

\[
\boxed{
\dim\ker\mathcal J_S^1=(k-1)d^k+b^k,
\qquad
\dim\operatorname{im}\mathcal J_S^1=2d^k-b^k.
}
\]

The proof of the image dimension is the exact sequence

\[
0\longrightarrow\epsilon\,\operatorname{im}\partial_N
\longrightarrow\operatorname{im}\mathcal J_S^1
\longrightarrow E\longrightarrow0.
\]

This also describes the precise quotient of the first thickening seen by the sum derivative. It does
not silently identify that quotient with the entire first thickening.

\paragraph{Retaining every first derivative is a different
observation}\label{conormal_tower_pasted:retaining-every-first-derivative-is-a-different-observation}

If all coordinate derivatives are retained, the map is

\[
\mathcal J_{\mathrm{all}}^1:
E^{[2]}\longrightarrow E\oplus E^k,
\qquad
[F]\longmapsto
\left([F],([\partial_iF])_i\right),
\]

where \(E^k\) is a square-zero ideal in the target.

On \(I/I^2\cong E^k\), its derivative component is

\[
(a_i)\longmapsto(c_i a_i)_i.
\]

Each multiplication map \(M_{c_i}\) has kernel dimension \(b\,d^{k-1}\). Hence

\[
\boxed{
\dim\ker\mathcal J_{\mathrm{all}}^1
=
k\,b\,d^{k-1}.
}
\]

In the squarefree case this full first-jet map is an algebra isomorphism onto the square-zero
extension \(E\oplus E^k\). With repeated roots it still loses the explicitly computed
Jacobian-annihilator directions, which remain in the higher source layers.

The split lifts retain their original meaning throughout:

\[
G(\pi)(n^\bullet)=e_E,
\qquad
G(\mathcal J_S^1)(n^\bullet)
=
(\epsilon\,\partial_Nn)^\bullet.
\]

External absence maps to external absence. This is a calculation of a retained relation before
quotienting, not a rule that differentiates \(e\) into a nonzero scalar. The attachment makes that
distinction explicitly.

\subsubsection{3. This calculation restricts to the symmetric tensor
construction}\label{conormal_tower_pasted:this-calculation-restricts-to-the-symmetric-tensor-construction}

The preceding maps commute with simultaneous permutation of variables and conormal indices. Let

\[
N=I/I^2\cong E^k,\qquad
T=\mathbb C[s]/(h,h').
\]

The cokernel sequence is equivariant:

\[
N\xrightarrow{\partial_N}E\longrightarrow T^{\otimes k}\longrightarrow0.
\]

Taking \(S_k\)-invariants remains exact: given an invariant element in an image, average any
preimage with the literal factor \(1/k!\). Thus

\[
\boxed{
\operatorname{coker}\left(
\partial_N:N^{S_k}\to E^{S_k}
\right)
\cong
(T^{\otimes k})^{S_k}.
}
\]

The invariant conormal domain has the explicit description

\[
\boxed{
N^{S_k}
\cong
A_h\otimes
(A_h^{\otimes(k-1)})^{S_{k-1}}.
}
\]

An invariant tuple is determined by its first component, which is invariant under permutations
fixing the first index; the other components are its prescribed translates. This constructs both
directions of the identification.

Consequently,

\[
\dim N^{S_k}
=
d\binom{d+k-2}{k-1},
\]

\[
\boxed{
\operatorname{rank}\partial_N^{S_k}
=
\binom{d+k-1}{k}
-
\binom{b+k-1}{k},
}
\]

and

\[
\boxed{
\dim\ker\partial_N^{S_k}
=
d\binom{d+k-2}{k-1}
-\binom{d+k-1}{k}
+\binom{b+k-1}{k}.
}
\]

For \(b=0\), the final binomial is zero.

For \(h=s^2\) and \(k=2\), the symmetric conormal domain has dimension four, the symmetric
arithmetic target has dimension three, and the restricted derivative has rank two. Its kernel still
has dimension two.

These are ordinary variable-permutation invariants of the relation algebra. To place them in the
original cochain construction, use the source's \textbf{signed chain idempotent}

\[
\mathsf S_k=\frac1{k!}\sum_{\pi\in S_k}\operatorname{sgn}(\pi)T_\pi.
\]

Its two top-degree signs produce precisely the ordinary symmetric tensors used here. Replacing it by
the unsigned chain average would change the selected summand.

This gives the smaller cohomological calculation its own exact derivative kernel; selecting the
symmetric summand does not amount to presuming that the remaining relative information vanishes.

\subsubsection{4. The arithmetic unit and trace remain
attached}\label{conormal_tower_pasted:the-arithmetic-unit-and-trace-remain-attached}

The new note retains

\[
U=\prod_i(g/h)(s_i)
\]

and its Taylor class \(\widehat U\in E^{[2]}\). The relative Leibniz rule proves

\[
\boxed{
\delta_1(\widehat Ux)
=
U\,\delta_1(x)+\delta_1(\widehat U)\pi(x).
}
\]

On a conormal class \(n\), the second term vanishes:

\[
\boxed{
\delta_1(\widehat Un)
=
U\,\partial_Nn.
}
\]

Because \(U\) is a unit, this arithmetic multiplication changes neither the kernel nor the ranks
calculated above. It is nevertheless retained in the actual map, rather than assigned the value one.
This is the source's specified unit correction.

There is also an important distinction between the \textbf{sum-Jacobian response} and the
\textbf{product Jacobian used by the full product trace}.

Multiplication by

\[
\prod_i h'(s_i)
\]

on \(E=A_h^{\otimes k}\) has rank

\[
\boxed{q^k,}
\]

where \(q\) is the number of distinct roots of \(h\). Its kernel has dimension

\[
d^k-q^k
\]

and is the nilradical of \(E\).

\textbf{Proof.} At a root \(\rho\) of multiplicity \(m\), write

\[
h(\rho+z)=z^m u_\rho(z),\qquad u_\rho(0)\ne0.
\]

Modulo \(z^m\),

\[
h'(\rho+z)=m\,u_\rho(0)z^{m-1}.
\]

Multiplication by this element has one-dimensional image and kernel \((z)\). At a \(k\)-tuple of
roots, the product Jacobian similarly sends the constant coefficient to a nonzero multiple of

\[
z_1^{m_1-1}\cdots z_k^{m_k-1}
\]

and annihilates the maximal ideal. There is one image dimension per root tuple, giving \(q^k\).

This is not the rank \(d^k-b^k\) of the contracted conormal row. The two operators use the same
derivative factors but perform different observations.

The one-variable trace identity remains exact. Let

\[
\lambda_h([P])
=
\sum_{\rho}\operatorname{Res}_{s=\rho}\frac{P(s)}{h(s)}\,ds.
\]

Then

\[
\boxed{
\lambda_h(h'P)
=
\operatorname{Tr}(M_P\mid A_h).
}
\]

Locally, \(h'/h=m_\rho/(s-\rho)+u_\rho'/u_\rho\), so its residue against \(P\) is
\(m_\rho P(\rho)\). Multiplication by \(P\) on the local algebra has that same trace because its
nilpotent terms have zero diagonal.

Writing \(\upsilon_h=j_h(g/h)\) and \(\varepsilon_h=\upsilon_h^{-1}\), the actual arithmetic pairing
gives

\[
\begin{aligned}
\mathscr R_Z(f,j_h(g')P)
&=\lambda_h\!\left(
\varepsilon_h f^\dagger\upsilon_h h'P
\right)\\
&=\lambda_h(f^\dagger h'P)\\
&=\operatorname{Tr}(M_{f^\dagger P}\mid A_h).
\end{aligned}
\]

Both unit factors are present before cancellation. This is the algebraic form of the attachment's
observation of

\[
(\log x)\Theta(P(D)\phi_*),
\]

with the original \(g=2\xi\), reflection, and multiplicities retained.

\subsubsection{5. Lean continuation against the integrated
library}\label{conormal_tower_pasted:lean-continuation-against-the-integrated-library}

The merged Codex source already supplies

\begin{verbatim}
SplitZero.RelationLayer.derivative
\end{verbatim}

as the quotient map induced by a linear map that carries the source relations into the target
relations. I read its implementation and reused that interface below rather than introducing another
quotient.

The following proposed \nolinkurl{SplitZeroConormalTower.lean} proves the ideal-power containment from
the derivation law, uses it to construct every descended derivative, proves compatibility between
levels, and computes the derivative of an original relation. Its proof scripts are complete, but
\textbf{compilation has not been checked}.

\begin{verbatim}
import SplitZeroRelationLayerIntegration
import Mathlib.RingTheory.Derivation.Basic
import Mathlib.RingTheory.Ideal.Operations

/-!
Derivative maps between the original ideal-power quotients.

This draft reuses RelationLayer.derivative. It does not assume that
the ideal is derivative-stable or that the packet is squarefree.
-/

noncomputable section

namespace SplitZero.ConormalTower

variable {R A : Type*}
  [CommRing R] [CommRing A] [Algebra R A]

/-- The original ideal-power relation, viewed over the coefficient ring. -/
def level (I : Ideal A) (r : ℕ) : Submodule R A :=
  (I ^ r).restrictScalars R

/-- A derivation lowers ideal order by at most one. -/
theorem deriv_mem_pow
    (D : Derivation R A A) (I : Ideal A) (r : ℕ)
    {x : A} (hx : x ∈ I ^ (r + 1)) :
    D x ∈ I ^ r := by
  induction r generalizing x with
  | zero =>
      rw [Submodule.pow_zero, Ideal.one_eq_top]
      exact Submodule.mem_top
  | succ r ih =>
      rw [Submodule.pow_succ] at hx
      refine Submodule.mul_induction_on hx ?_ ?_
      · intro a ha b hb
        rw [D.leibniz]
        change a * D b + b * D a ∈ I ^ (r + 1)
        have hd : D a * b ∈ I ^ (r + 1) := by
          rw [Submodule.pow_succ]
          exact Ideal.mul_mem_mul (ih ha) hb
        exact (I ^ (r + 1)).add_mem
          ((I ^ (r + 1)).mul_mem_right (D b) ha)
          (by simpa only [mul_comm] using hd)
      · intro a b ha hb
        rw [map_add]
        exact (I ^ (r + 1)).add_mem ha hb

/-- The derivative on A/I^(r+1), with values in A/I^r. -/
def descended
    (D : Derivation R A A) (I : Ideal A) (r : ℕ) :
    (A ⧸ level (R := R) I (r + 1)) →ₗ[R]
      (A ⧸ level (R := R) I r) :=
  SplitZero.RelationLayer.derivative
    (level I (r + 1)) (level I r) D.toLinearMap
    (fun _ hx => deriv_mem_pow D I r hx)

@[simp]
theorem descended_mk
    (D : Derivation R A A) (I : Ideal A) (r : ℕ) (x : A) :
    descended D I r ((level I (r + 1)).mkQ x) =
      (level I r).mkQ (D x) :=
  rfl

/-- Forget one ideal-adic layer without differentiating. -/
def projection (I : Ideal A) (r : ℕ) :
    (A ⧸ level (R := R) I (r + 1)) →ₗ[R]
      (A ⧸ level (R := R) I r) :=
  SplitZero.RelationLayer.transport
    (level I (r + 1)) (level I r)
    (fun _ hx => Ideal.pow_le_pow_right (Nat.le_succ r) hx)

@[simp]
theorem projection_mk (I : Ideal A) (r : ℕ) (x : A) :
    projection (R := R) I r ((level I (r + 1)).mkQ x) =
      (level I r).mkQ x :=
  rfl

/-- Derivative and quotient transition form the actual commuting square. -/
theorem tower_square
    (D : Derivation R A A) (I : Ideal A) (r : ℕ) :
    (projection I r).comp (descended D I (r + 1)) =
      (descended D I r).comp (projection I (r + 1)) := by
  apply LinearMap.ext
  intro z
  obtain ⟨x, rfl⟩ :=
    Submodule.Quotient.mk_surjective (level I (r + 2)) z
  rfl

/-- On a retained relation z*a, only a*D(z) survives modulo I. -/
theorem conormal_product
    (D : Derivation R A A) (I : Ideal A)
    (z a : A) (hz : z ∈ I) :
    descended D I 1 ((level I 2).mkQ (z * a)) =
      (level I 1).mkQ (a * D z) := by
  rw [descended_mk, D.leibniz]
  change
    (level I 1).mkQ (z * D a + a * D z) =
      (level I 1).mkQ (a * D z)
  rw [map_add]
  have hzero : (level (R := R) I 1).mkQ (z * D a) = 0 := by
    apply (Submodule.Quotient.mk_eq_zero _).mpr
    change z * D a ∈ I ^ 1
    rw [Submodule.pow_one]
    exact I.mul_mem_right (D a) hz
  rw [hzero, zero_add]

end SplitZero.ConormalTower
\end{verbatim}

The derivation and ideal-product APIs used here were checked against the repository's existing
Mathlib revision, rather than assumed from a different release. That source inspection does not
substitute for running Lean.

\subsubsection{Verification and the remaining quantitative
connection}\label{conormal_tower_pasted:verification-and-the-remaining-quantitative-connection}

The fresh Wolfram calculations passed \textbf{161 exact checks}: conormal ranks and invariant
restrictions, full first-jet ranks, ideal-square annihilation, residue traces, product-Jacobian
ranks, higher derivative identities, and retained-unit trace identities. The fixtures include
squarefree polynomials, repeated roots, mixed multiplicities, and tensor degrees through three. An
initial variable-binding error in my checker was corrected before the successful runs. These are
finite algebra checks, not a new Lean certificate or evaluations of arithmetic zero locations.

The new result does not identify the Hilbert-space normal map \(N_u\) with the algebraic conormal
module \(I/I^2\). The attachment's filtered graph supplies the stated route from the full derivative
to arithmetic jets. Its exact scalar identity

\[
\mathcal I_{h,k}
+4\int\|n_k(u)\|^2\,du
=
\frac{\mu_h^{k-1}}{k}\mathcal I_h
\]

retains the Hilbert normal contribution; the conormal calculations above now identify the algebraic
kernel encountered by the subsequent sum-direction observation. Those two pieces must remain
separately typed when estimating the original weight-control form.

\textbf{The continuation now provides an exact first-jet algebra map, its full retained kernel, the
corresponding symmetric restriction, and all-order quotient derivatives.} Codex can build those on
the merged quotient library without recreating the scalar, homology, or boundary-control
implementations.

\paragraph{Reference}\label{conormal_tower_pasted:reference}

The Stacks Project Authors. (n.d.). \emph{The naive cotangent complex} (Tag 00S0). \emph{The Stacks
Project}. (\href{https://stacks.math.columbia.edu/tag/00S0}{The Stacks Project})
